\section{introduction}

The equivalent-layer technique was introduced during the 1960s and has become a versatile
tool for processing and interpreting potential-field data in applied geophysics . 
The technique can be used to accomplish various data processing steps, such as interpolation 
\citep[e.g.,][]{cordell_1992,mendonca-silva_1994, barnes2012, soler-uieda2021}, 
upward (or downward) continuation 
\citep[e.g.,][]{henderson-zietz1949, emilia1973, bhattacharyya-chan1977,
hansen-miyazaki_1984, li-oldenburg_2010, siqueira_etal_2017, takahashi-etal2020},
gradient data processing 
\citep[e.g.,][]{barnes-lumley_2011, barnes2014, martinez_li_2016}, 
merging multiple data sets \citep[e.g.,][]{boggs-dransfield2004, lane2004}, 
estimating density models from gravity data \cite[e.g.,][]{mendonca-2020}, 
computing the amplitude of the anomalous magnetic field 
\citep[e.g.,][]{li_li_2014, sun_etal-2019, melo-etal2021, hidalgo-gato-etal2021}, 
transforming total-field anomalies into field components \citep[e.g.,][]{li-etal2021}, 
as well as estimating the magnetization direction and computing the reduction to the pole
of magnetic data \citep[e.g.,][]{emilia1973, silva_1986, leao-silva_1989, 
guspi-novara_2009, oliveirajr-etal_2013, li_etal_2014, oliveirajr_etal_2015, 
ribeiro_filho_etal2020, reis_etal2020}.

Methods for processing/interpreting potential-field data 
via equivalent-layer technique usually follow the pioneer approach presented
by \citet{dampney1969} and rely on the numerical solution of the 
\textit{upward continuation integral} \citep[e.g.,][]{skeels1947, peters1949, henderson-zietz1949,
dean1958, roy1962, bhattacharyya1967, henderson1970}, which is the solution of the
\textit{Dirichlet’s problem} or the \textit{first boundary value problem of potential
theory} on a plane \cite[][ p. 236]{kellogg1967}. 
This numerical solution is analogous to the density (or mass) distribution on a
fictitious discrete layer of monopoles (points of mass) that are placed on a plane above 
the true sources and reproduce the discrete potential-field data set.
This layer is usually called \textit{equivalent layer} and the monopoles (or points of mass)
\textit{equivalent sources}.
By using the estimated physical property (density or mass) distribution of the 
equivalent sources, the interpreter may compute, for example, the upward/downward
continuation or interpolation.

Note that, when defined in this way, the equivalent layer is grounded on a solid 
theoretical background and represents the solution of a boundary value problem of 
potential theory. 
This theoretical background guarantees the existence of a planar equivalent layer of
monopoles that is located above the true sources, has a physical property distribution
defined by a scaled version of the observed potential field on the layer, 
reproduces the potential-field data at the observation points and 
allows computing some potential field transformations at points with no data above the
layer.

It is important to stress that the equivalent layer defined in space domain as the solution
of Dirichlet's problem on a plane is valid for the case in 
which the layer does not cross the sources. 
This condition, however, cannot be assured in practical situations with field data.
It is also important noting that this space domain approach does not guarantee that 
the same equivalent layer of monopoles reproducing conventional gravity data also 
reproduces the corresponding gravity gradient tensor components due the same sources.
Notwithstanding this lack of theoretical background, the equivalent-layer technique 
has been successfully used for jointly processing gravity-gradient 
data \citep[e.g.,][]{li2001, barnes-lumley_2011, martinez_li_2016} or converting gravity 
disturbance into gradient tensor components \citep[e.g.,][]{boggs-dransfield2004,  piauilino_etal2019_seg, li_etal2021}, for example.

In the geophysical literature, it has long been assumed by principle that the magnetic
induction field produced by 3D sources can be reproduced by a fictitious and 
purely mathematical planar layer of dipoles \citep[e.g.,][]{emilia1973, 
bhattacharyya-chan1977, silva_1986, oliveirajr-etal_2013, li_etal_2014, 
oliveirajr_etal_2015, ribeiro_filho_etal2020, li-etal2021}. 
Few works have investigated the theoretical aspects related to the existence of
a planar equivalent layer of dipoles or its relationship with a boundary value
problem of potential theory.
By following different approaches in the wavenumber domain, 
\citet{pedersen1991} and \citet{li_etal_2014} proved the existence of a
planar equivalent layer of dipoles with all-positive
magnetic moment distribution that exactly reproduces the approximated total-field anomaly
\cite[e.g., ][ p. 179]{blakely1996} caused by sources having a purely induced magnetization
in the vertical direction.
Their approach, however, is valid only for the case in which 
the total-field anomaly data rely on a plane.
\citet{reis_etal2020} generalized the works of \citet{pedersen1991} and \citet{li_etal_2014}
to prove that this planar equivalent layer exists for the case in which the dipoles have 
the same arbitrary uniform magnetization direction of the sources, regardless of whether 
the sources magnetization is purely induced or not.

Here, we present a more general discussion on the existence of planar equivalent layers
that exactly reproduce the gravitational and magnetic induction fields due to 3D sources,
even for the case in which the layer crosses them.
We limit ourselves to the case in which the magnetic
sources have a total-magnetization with uniform direction, but spatially-variable intensity.
We deduce the space-domain solutions of the Dirichlet's and Neumann's problems  
on a planar surface that crosses the sources. 
Then we use these solutions to prove the existence of a continuous planar equivalent layer
of monopoles that simultaneously reproduces the gravitational potential, attraction and
tensor components due to arbitrary gravity sources.
We also expand the work of \citet{reis_etal2020} to prove, in space domain, 
the existence of a continuous planar equivalent layer of dipoles that is all positive,
has the same constant total-magnetization direction of the true 3D sources and 
simultaneously reproduces the magnetic scalar potential and magnetic induction
field components along arbitrary constant directions.
Besides, we show that this planar equivalent layer of dipoles also reproduces the 
special magnetic anomaly that would be yielded by the true sources if they have
a purely induced magnetization with vertical direction.
This theoretical result has been implicitly assumed to be true, but never proved 
neither in the pioneer work of \citet{silva_1986}, nor in the succeeding 
methods \citep[e.g.,][]{leao-silva_1989, oliveirajr-etal_2013, li_etal_2014, reis_etal2020}
proposed to compute the reduction to the pole (RTP) in space domain.

The theoretical developments presented here open the door for future 
investigations in space domain about,
for example, the instability of RTP in space domain and its generalization 
to the case in which the sources have a nonuniform total-magnetization direction,
the existence of a closed-form solution for the case in which the planar equivalent layer 
of dipoles has a constant magnetization direction different from that of the true sources
and also indicate a research direction to understand the relationship between the planar
equivalent layer and the true 3D sources.